\section{Geodesic analysis details}
\label{app:geodesic-details}

\paragraph{Pairing and alignments.}
We draw 1000 molecules from the dataset and group them by node count. For each group we sample
pairs uniformly and align them using node matching to define a coupling between endpoints.

\paragraph{Path optimization.}
We represent a continuous path in the embedding space and optimize the energy-weighted length
objective in \Cref{sec:experiments}, including the path-energy term. The linear baseline uses the
same endpoint pairing and interpolates between the coupled graphs.

\paragraph{Spline parameterization.}
We implement geodesic paths using a Bezier spline over learnable control-point logits for nodes and
edges. The spline lives in the probability simplex (node/edge categorical distributions), and
discrete graphs are sampled only for visualization and metric evaluation. The optimization minimizes
the relative length ratio $L_{\mathrm{spline}}/L_{\mathrm{linear}} - 1$, where the segment lengths are
weighted by $\exp(-\beta V_\theta(x_\tau))$ along the path.

\paragraph{Discrete evaluation.}
We discretize each continuous path at evenly spaced points in normalized arc length. At each point
we sample categorical graphs from the path distribution to compute validity and energy statistics.
We report the length ratio $L_{\mathrm{geo}}/L_{\mathrm{lin}}$ and the correlation between energy
and validity along the path.
